\documentclass{article}
\usepackage[T1]{fontenc}
\usepackage[utf8]{inputenc}

% Linespread command allows you to change line spacing for the entire document
\linespread{1.18}

% Tweak page margins
\addtolength{\oddsidemargin}{-.875in}
\addtolength{\evensidemargin}{-.875in}
\addtolength{\textwidth}{1.75in}

\addtolength{\topmargin}{-.875in}
\addtolength{\textheight}{1.75in}

\usepackage{natbib}
\usepackage{hyperref}
\usepackage{xcolor}
\usepackage{xspace}
\usepackage{fancyhdr}
\hypersetup{
    colorlinks,
    linkcolor={red!50!black},
    citecolor={blue!50!black},
    urlcolor={blue!80!black}
}

\newcommand{\HRule}{\rule{\linewidth}{0.5mm}}
\newcommand{\Hrule}{\rule{\linewidth}{0.3mm}}

% Project specific macros
\newcommand{\graphite}{GRAPHITE\xspace}
\newcommand{\wave}{WAVE\xspace}

% School specific macros
\newcommand{\schoolShort}{Stanford\xspace}
\newcommand{\school}{Stanford\xspace}
\newcommand{\schoolLong}{\xspace}

\newcommand{\profOne}{Prof. Matei Zaharia\xspace}
\newcommand{\profTwo}{Prof. Christos Kozyrakis\xspace}
\newcommand{\profThree}{Prof. Jure Leskovec\xspace}

% Creates header for each page
\usepackage{fancyhdr}
\pagestyle{fancy}
\fancyhf{}
\fancyhead[LE,RO]{\header\hskip\linepagesep\vfootline\thepage}
\newskip\linepagesep \linepagesep 5pt\relax
\def\vfootline{%
    \begingroup
    	\rule[-10pt]{0.75pt}{25pt}
    \endgroup
}
\def\header{%
	\begin{minipage}[]{120pt}
		\hfill Arunima Chowdhury 		% Applicant Name
    	\par \hfill 				% Formatting boilerplate
    	CS, PhD, Fall 2027 			% Area, Program, Cycle, Year
    \end{minipage}
}
\fancyhead[RE,LO]{Statement of Purpose | \schoolLong}
\renewcommand\headrulewidth{0pt}

\begin{document}

%% Why do you wish to attend graduate school? What would you like to study? Keep it broad, details come-in later
\textbf{Introduction:} My research interests center on making foundation models more reliable by studying how they reason and improving how we supervise and optimize them. My interest was piqued when I encountered a gap in my medical image reasoning research, where visual perception and clinical inference errors need distinct feedback, yet current training methods rarely provide it. In domains where correctness is easy to check, such as math and code, models can learn to reason in steps when trained with verifiable rewards. Additional computation at test time can further improve performance. In other settings, learned reward models provide feedback, but their judgments can be unreliable and models can learn to exploit them. In my PhD, I want to develop step-level supervision that helps identify where reasoning errors originate and study how data selection and optimization can make this feedback reliable enough to improve training and test-time answer ranking.



%% Describe 2-3 past projects that might be relevant to your research interests. (10-12 lines per project)

% PROJECT 1: med prm

\textbf{Research Experience:} I first encountered this challenge during my master's research on multimodal medical reasoning. While studying process reward models, I became interested in a key limitation: they can identify incorrect reasoning steps, but often provide little insight into what caused the error. For medical image analysis, a model must locate the correct region, recognize what it contains, and reason from it to a diagnosis. A model might misinterpret an image or draw an incorrect clinical conclusion from accurately identified visual evidence. Although both lead to incorrect reasoning, they represent different types of errors that require different forms of supervision.

To address this, I am leading a research on a typed process reward model that judges each reasoning step on two separate axes: whether its visual claims are correct, and whether its clinical inference holds given those claims. To build supervision without manually labeling each reasoning step, I sample reasoning chains from Qwen2.5-VL-7B, parse them into visual claims and inferences, check visual claims against available expert annotations, and evaluate clinical inferences against retrieved medical knowledge. If an inference relies on a visual claim that turns out to be wrong, I label it as dependent rather than counting it as a separate error. After testing by replacing the real image with a blank one, the accuracy drops by \textbf{16\%}, suggesting the model utilizes the image for its assessments. I plan to use the trained model both for Best-of-N selection among candidate reasoning chains and as a step-level reward in GRPO to train models that reason from the image rather than from language priors. Designing this PRM taught me that how we define an error shapes what a model can learn to correct. It made me pay closer attention to which judgments we ask a reward model to make before using its feedback to guide training.

% PROJECT 2: bangla

I encountered a related challenge in an earlier project on Bangla speech emotion recognition, where labeled emotional speech is scarce. After fine-tuning Whisper and HuBERT on a small labeled dataset, I found that performance varied sharply across emotion classes, especially those with few examples. To make use of a larger unlabeled corpus, I implemented iterative pseudo-labeling, where I added the models' high-confidence predictions to subsequent training rounds. But confidence alone was unreliable: confidently wrong predictions reinforced existing class biases, particularly between acoustically similar emotions. To tackle this problem, I combined confidence filtering with class-balanced selection. This raised macro-F1 from 61.8\% to 68.4\%, with the largest gains in underrepresented classes. The experience taught me that adding more data is not enough if the supervision itself is unreliable. What matters is whether we can identify and control errors in the labels we learn from, a concern I now revisit when constructing step-level labels for reasoning.

% PROJECT 3: Dance Motion Feedback System

My undergraduate research at Purdue gave me an early appreciation for how much the quality of model-generated feedback depends on the evidence available to it. I spent two years building an automated feedback system for modern dance that compares student performances against instructor references and generates personalized corrective guidance. The project started with 3D motion tracking for live theatre, where I needed fast, reliable pose estimation for multiple dancers. YOLO-Pose with stereo depth sensing gave me that; roughly 3x faster and 70\% more accurate than TensorFlow PoseNet. I classified movements across 19 categories using a fine-tuned SkateFormer and aligned student and instructor performances using FastDTW to measure joint-level deviations. To turn these measurements into actionable coaching, I used Gemini as a feedback generator. I provided it with structured, joint-level movement deviations and prompted it to generate corrective feedback tied to those measurements, rather than relying on the model to identify movement errors on its own. The work received the Best Research Talk Award at Purdue's Undergraduate Research Conference.

%% Non-research accomplishments (e.g. Grades, Academic Service, Work experience) (10-12 lines)
\textbf{Diversity and Service: }Beyond research, I have sought opportunities to develop the technical judgment, communication, and leadership skills needed for collaborative academic work. At Purdue, I was recognized on the Dean's List for academic performance. During my master's, I served as a senior grader evaluating student teams in project-based software engineering courses. Reviewing their designs, implementations, and individual contributions taught me to assess technical work critically, explain weaknesses constructively, and remain fair when contributions were uneven.

I developed my leadership skills at Purdue's Data Mine program, where I led nine student developers in a collaboration with Merck. Coordinating milestones and communicating with industry collaborators taught me how to translate broad research objectives into concrete tasks while keeping a team aligned. My internship at Robi Axiata gave me experience applying machine learning to real-world problems, where I developed forecasting models using thousands of subscriber records, achieving approximately 12\% MAPE. Beyond technical work, I served as Vice President of the Bangladesh Student Association, organizing cultural programs, fostering community engagement, and strengthening connections between Bangladeshi students and the broader campus community. Together, these experiences taught me how to balance independent judgment with teamwork, communicate across different backgrounds, and take responsibility for work that involves others.

%% Why this school? List professors you would like to work with and why? (10-12 Lines)

At George Mason, I want to study how reliable feedback can improve model training and answer selection at test time. I am really interested in working with Prof. Ziwei Zhu, whose research on credit assignment aligns with my goal of developing better supervision for reasoning models. I want to explore if better supervision is achievable by separating the source of an error from its later effects. I am also drawn to Prof. Ziyu Yao's work on multimodal reasoning and interpretability, which complements my interest in understanding why models make errors. I am excited about using interpretability to understand what internal mechanisms a model relies on when it produces confident but incorrect reasoning. With Prof. Mingrui Liu, I would like to explore how data selection and optimization can improve the way models learn from feedback. I am also eager to work with Prof. Kenton Murray, whose research on multilingual learning and preference optimization connects to an interest I developed through my Bangla speech research in learning from limited labeled data.


%% Summary (3-4 Lines)

These skills, combined with my research experience, have prepared me to take on the challenges of doctoral work. I want to pursue a PhD to go deeper into the questions my work has raised: how to build supervision that models can reliably learn from, especially when correctness is hard to verify. Long term, I want to build a career as a research scientist, where I can pursue these questions independently, lead collaborative research projects, and develop methods that make foundation models more dependable in the settings where they matter most. I see doctoral training as an important step toward becoming a researcher capable of identifying fundamental problems and developing practical ways to address them.

% TEMPLATE: Replace bracketed placeholders for each school.

%After completing my PhD, I plan to [CAREER GOAL — e.g., pursue a career as a research scientist working on X / continue in academia]. My experience with [BRIEF CALLBACK to your work] has shown me that [CORE BELIEF about your research area]. [SCHOOL]'s mentorship and resources will enable me to [FORWARD-LOOKING STATEMENT].

%Across these experiences, I have become increasingly interested in understanding not just whether a model produces the correct answer, but whether the feedback used to train and evaluate it can be trusted. I want to pursue a PhD to study how we can build more reliable supervision for foundation models, especially in tasks where correctness cannot be easily verified. I am particularly interested in developing reward models that can identify where reasoning goes wrong and exploring how these signals can improve reinforcement learning and inference-time reasoning. In the long term, I hope to build a career in AI research, where I can pursue these questions independently, lead collaborative research projects, and develop methods that make foundation models more dependable in real-world applications. I see doctoral training as an important step toward becoming a researcher capable of identifying fundamental problems and developing practical ways to address them.


% Add some blank space between text and references
\vspace{0.125in}

% References

% **NOTE**: There are better ways to manage citations in LaTeX, most notably using a bibTeX. I wanted to have greater control on how citations were spaced and formatted and therefore ended up hardcoding them here. Your mileage may wary!



\end{document}

% That's All Folks.

% Best of luck, you got this! :)
